\begin{abstract}
Several methods have been proposed and used to refine estimates of future climate change based on combined output from comprehensive climate models. While previously the so-called model democracy approach was used to combine model predictions, where every model is given equal weight, it is now widely accepted that using model weights that account for model performance and model independence yields better results. However, most existing approaches rely, implicitly or explicitly, on a similar statistical basis, while describing things in different ways. Here we aim to build a common framework formulated in probabilistic Bayesian terms and distinguishing between approaches that are based on the performance of individual models (individual performance weighting) and approaches that are based on the performance of the weighted ensemble as a whole (ensemble performance weighting). Using simple constructed examples, we demonstrate that the ensemble performance weighting approach implicitly accounts for co-dependencies among models, and examine how methodological choices such as the prior distribution affect the results. We attempt to clearly define the characteristics of the task in general terms, and from there, arrive at useful guidelines and strategies to make robust climate projections with uncertainty estimates, as well as establishing a common basis for future developments with regard to ensemble model weighting in Earth system science.
\end{abstract}